% 数理物理1 数理物理1-12 問題1（092-bunsan。当時は 数理物理1-10 問題1・shou10-l1-bunsan） 解いた日 2026-10-07
%% kotae
\toi{1}{◎}{}
\begin{align*}
  \omega_\mu&=\sqrt{\frac{2K}{m}\cdot2\sin^{2}\frac{\pi\mu}{N}}\\
  &=2\sqrt{\frac{K}{m}}\sin\Bigl(\frac{\pi\mu}{N}\Bigr)\\
  &=2\sqrt{\frac{K}{m}}\sin\frac{k_\mu a}{2}
\end{align*}
\toi{2}{◎}{}
$\sin(k_\mu a/2)\simeq k_\mu a/2$ より，
\[ \omega_\mu=\underbrace{a\sqrt{\frac{K}{m}}}_{v_s}k_\mu,\qquad v_s=a\sqrt{\frac{K}{m}} \]
\toi{3}{△}{単位}
\begin{align*}
  v_s&=2.5\times10^{-10}\cdot\sqrt{\frac{15}{4.5\times10^{-26}}}\\
  &=2.5\cdot\frac{\sqrt{10}}{\sqrt{3}}\cdot10^{3}\\
  &=\ayamari{4.56\times10^{3}}{\qty{4.6e3}{\metre\per\second}}\\
  \omega_{\max}&=2\sqrt{\frac{15}{4.5\times10^{-26}}}\\
  &=\ayamari{3.65\times10^{13}}{\qty{3.7e13}{\per\second}}\\
  \nu_{\max}&=\frac{\omega_{\max}}{2\pi}=\ayamari{5.81\times10^{12}}{\qty{5.8e12}{\hertz}}
\end{align*}
\naze{数値は正しいが，3つとも単位が無く，指定の有効数字2桁にもなっていない。$\omega_{\max}$ は $\abs{\sin(k_\mu a/2)}=1$（$k_\mu a=\pi$，$\mu=N/2$）のときであることも一言添える。}
\toi{4}{×}{立式}
\begin{align*}
  \frac{\omega_\mu}{v_sk_\mu}&=\frac{2\sqrt{K/m}\,\sin(k_\mu a/2)}{a\sqrt{K/m}\,k_\mu}\\
  &=2\sin\frac{1}{2}=2\Bigl(\frac{1}{2}-\frac{1}{6}\Bigl(\frac{1}{2}\Bigr)^{3}\Bigr)\\
  &=1-\frac{1}{24}=\frac{23}{24}
\end{align*}
\[ \ayamari{2\sin x}{\frac{2\sin(x/2)}{x}}=\frac{99}{100}\qquad\therefore\ 2x-\frac{1}{3}x^{3}=\frac{99}{100} \]
\naze{比は $x=k_\mu a$ の関数として $2\sin(x/2)/x$ である。1行目では $k_\mu a=1$ を入れたので分母が 1 になって消えたが，一般の $x$ では分母が残る。正しくは $2\sin(x/2)/x\simeq1-x^{2}/24=0.99$ から $x^{2}=0.24$，$k_\mu a\simeq0.49$，$\lambda_\mu=2\pi a/x\simeq13a$。ここで止まっていて，$k_\mu a$ と $\lambda_\mu/a$ が出ていない。}
%% kizuki
\item 数値の答えには単位を付け，問題の指定（有効数字2桁）に丸めてから書く。$v_s$ は \unit{\metre\per\second}，$\omega_{\max}$ は \unit{\per\second}，$\nu_{\max}$ は \unit{\hertz} である。
\item 「この比が $0.99$ になる $x$」を求めるときは，値を代入する前の一般の式に戻って立式する。$x=k_\mu a$ とおくと比は $2\sin(x/2)/x$ で，$\sin(x/2)\simeq x/2-x^{3}/48$ を入れると $1-x^{2}/24$ になる。前半で $x=1$ を入れたときの $2\sin(1/2)$ をそのまま使うと分母の $x$ が消える。立てた式に $x=1$ を入れて前半の値 $0.96$ に戻るかで確かめる。
